\documentclass[11pt]{amsart}

\usepackage[T1]{fontenc}
\usepackage{lmodern}
\usepackage{microtype}
\usepackage{mathtools,amssymb,amsthm}
\usepackage{booktabs}
\usepackage[hidelinks]{hyperref}

\hypersetup{
  pdftitle={A Counterexample to the Hibi--Seyed Fakhari Conjecture on Join--Meet Ideals}
}

\newtheorem{theorem}{Theorem}

\title[A counterexample on join--meet ideals]
{A Counterexample to the Hibi--Seyed Fakhari Conjecture on Join--Meet Ideals}

\date{}

\subjclass[2020]{Primary 13D02; Secondary 05E40, 06B05}
\keywords{join--meet ideal, finite lattice, linear syzygies, graded Betti numbers}

\begin{document}

\begin{abstract}
Hibi and Seyed Fakhari conjectured that the join--meet ideal of every
nonmodular finite lattice is not linearly related. We refute their
conjecture with a $12$-element lattice. Let $L$ be the bounded
rank-three lattice whose atom--coatom incidence graph is the cycle
$C_{10}$. Over $\mathbb F_2$, and over every field of characteristic
zero,
\[
  \beta_{1,3}(I_{L,K})=149,
  \qquad
  \beta_{1,j}(I_{L,K})=0 \quad (j\ge 4).
\]
The vanishings in degrees four and five are certified by exact binary
rank computations. The range $j\ge6$ is not computed: it follows from
the regularity bound $\operatorname{reg}(I_{L,K})=4$, which is quoted
from Hibi and Seyed Fakhari.
\end{abstract}

\maketitle

\section{The counterexample}

Let $K$ be a field and $L$ a finite lattice. Set
$S_K=K[x_a:a\in L]$. Its join--meet ideal is
\[
  I_{L,K}
  =
  (x_ax_b-x_{a\wedge b}x_{a\vee b}:a,b\in L)
  \subseteq S_K.
\]
Write
$\beta_{i,j}(M)=\dim_K\operatorname{Tor}_i^{S_K}(M,K)_j$.
An ideal generated in degree two is \emph{linearly related} if
$\beta_{1,j}=0$ for every $j\ge 4$. Hibi and Seyed Fakhari conjectured
that $I_{L,K}$ is not linearly related whenever $L$ is nonmodular
\cite[Conjecture~2.4]{HSF}.

Define
\[
  L=
  \{\widehat 0,\widehat 1\}
  \cup
  \{a_i,b_i:i\in\mathbb Z/5\mathbb Z\}
\]
by taking $a_0,\ldots,a_4$ as the atoms and
$b_0,\ldots,b_4$ as the coatoms, with $b_i$ covering precisely
$a_i$ and $a_{i+1}$. Thus the incidence graph on the ten middle
vertices is
\[
  a_0-b_0-a_1-b_1-a_2-b_2-a_3-b_3-a_4-b_4-a_0.
\]
Indices are read modulo $5$.

\begin{theorem}\label{thm:main}
The lattice $L$ is nonmodular. If $K=\mathbb F_2$, or if $K$ has
characteristic zero, then
\[
  \beta_{1,3}(I_{L,K})=149,
  \qquad
  \beta_{1,j}(I_{L,K})=0
  \quad\text{for every }j\ge4.
\]
Consequently, Conjecture~2.4 of Hibi and Seyed Fakhari is false.
\end{theorem}

\begin{proof}
Two distinct atoms have a common coatom exactly when they are
consecutive, in which case that coatom is unique; otherwise their join
is $\widehat 1$. The dual statement holds for coatoms, and the mixed
meets and joins are immediate. Hence the displayed incidence relations
define a lattice.

The lattice is not modular. Indeed, $a_0\le b_0$, while
$a_2\wedge b_0=\widehat0$ and $a_0\vee a_2=\widehat1$. Therefore
\[
  a_0\vee(a_2\wedge b_0)=a_0
  \neq
  b_0=(a_0\vee a_2)\wedge b_0.
\]

Write
\[
  y=x_{\widehat0},\qquad z=x_{\widehat1},\qquad
  A_i=x_{a_i},\qquad B_i=x_{b_i}.
\]
If $a\le b$ then $a\wedge b=a$ and $a\vee b=b$, so a comparable pair
contributes the zero binomial. Of the $\binom{12}{2}=66$ pairs of
distinct elements, $31$ are comparable and drop out, and the generating
set of $I_{L,K}$ is indexed by the remaining $35$ incomparable pairs:
ten pairs of atoms, ten pairs of coatoms, and fifteen incomparable
atom--coatom pairs. The corresponding $35$ quadrics consist of
\[
  A_iA_{i+1}-yB_i,
  \qquad
  B_iB_{i+1}-A_{i+1}z
  \qquad(i\in\mathbb Z/5\mathbb Z),
\]
one for each of the five pairs of consecutive atoms and each of the five
pairs of consecutive coatoms, together with twenty-five binomials
$UV-yz$. Here $U$ and $V$ run over the pairs of distinct middle
variables whose lattice elements are incomparable with meet $\widehat0$
and join $\widehat1$, namely the five pairs of nonconsecutive atoms, the
five pairs of nonconsecutive coatoms, and the fifteen incomparable
atom--coatom pairs. The first monomials of the $35$ quadrics are
pairwise distinct and none of them occurs as a second monomial, so the
$35$ quadrics are linearly independent over $K$; since they span
$(I_{L,K})_2$ and $(I_{L,K})_1=0$, they are a $K$-basis of
$(I_{L,K})_2$ and hence a minimal system of generators of $I_{L,K}$.

Let $\mathcal G$ be this set of generators, set
\[
  F=\bigoplus_{g\in\mathcal G}S(-2)e_g,
  \qquad
  \varphi:F\longrightarrow I_{L,K},
  \qquad
  e_g\longmapsto g,
\]
where $S=S_K$, and write $Z=\ker\varphi$. The degree-two
independence gives $Z_2=0$.

For $d\ge2$, define a multigraph $\Gamma_d$ as follows. For every
$g=m_g-n_g\in\mathcal G$ and every monomial $u$ of degree $d-2$, add
an edge $(g,u)$ joining the degree-$d$ monomials $um_g$ and $un_g$.
The vertices are the monomials incident with at least one edge. After
composing $\varphi_d$ with the inclusion $(I_{L,K})_d\hookrightarrow
S_d$ and using the monomial basis of $S_d$, its matrix is an oriented
incidence matrix of $\Gamma_d$, with zero rows omitted. Thus, writing
$E_d,V_d,c_d$ for the numbers of edges, vertices, and connected
components, the graph-incidence rank formula gives
\[
  \dim_K Z_d=E_d-V_d+c_d.
\]
This dimension is independent of $K$.

Let $\mathfrak m=S_+$. Since
$(\mathfrak mZ)_d=S_1Z_{d-1}$,
\[
  \beta_{1,d}(I_{L,K})
  =
  \dim_K Z_d-\operatorname{rank}_K(\mu_d),
  \qquad
  \mu_d:S_1\otimes_K Z_{d-1}\longrightarrow Z_d.
\]
We now work over $\mathbb F_2$. Choose a spanning forest in each
$\Gamma_d$. Its non-tree edges determine the fundamental-cycle basis
of $Z_d$. In this basis, the coordinate of a cycle at a non-tree edge
is its incidence with that edge. Multiplication by a variable sends
an edge $(g,u)$ of $\Gamma_{d-1}$ to the edge $(g,x_ru)$ of
$\Gamma_d$, so the matrices of $\mu_4$ and $\mu_5$ are obtained
exactly by bit-vector arithmetic.

The graph data are
\[
\begin{array}{c@{\qquad}r@{\qquad}r@{\qquad}r@{\qquad}r}
\toprule
 d & E_d & V_d & c_d & \dim Z_d\\
\midrule
 3 &   420 &  312 &  41 &  149\\
 4 &  2730 & 1293 &  91 & 1528\\
 5 & 12740 & 4276 & 161 & 8625\\
\bottomrule
\end{array}
\]
and exact Gaussian elimination over $\mathbb F_2$ gives
\[
  \operatorname{rank}_{\mathbb F_2}(\mu_4)=1528
  \quad\text{from a }1788\times1528\text{ matrix},
\]
\[
  \operatorname{rank}_{\mathbb F_2}(\mu_5)=8625
  \quad\text{from a }18336\times8625\text{ matrix}.
\]
In both displays the matrix is written with its rows indexed by the basis
$\{x_r\otimes\zeta\}$ of $S_1\otimes_KZ_{d-1}$, so that
$1788=12\cdot149$ and $18336=12\cdot1528$, and with its columns indexed
by the fundamental-cycle basis of the target $Z_d$. In each case the
stated rank is therefore the full target rank. Consequently,
\[
  \beta_{1,3}(I_{L,\mathbb F_2})=149,
  \qquad
  \beta_{1,4}(I_{L,\mathbb F_2})
  =\beta_{1,5}(I_{L,\mathbb F_2})=0.
\]
We now check the hypotheses under which \cite[Lemma~2.8]{HSF} is stated,
namely those of \cite[Lemma~2.6]{HSF}. The lattice $L$ is finite and
pure of rank three: every maximal chain has the form
$\widehat0<a_i<b_j<\widehat1$ with $a_i\le b_j$, and so has length
three. Its Hasse diagram is connected, and $|L|=12\ge7$. Finally the
graph $G_L:=L\setminus\{\widehat0,\widehat1\}$ of
\cite[Section~2]{HSF} is the atom--coatom incidence graph displayed
above, namely the cycle $C_{2n}$ with $2n=10\ge8$, which is connected.
All hypotheses hold, so
\cite[Lemma~2.8]{HSF} applies and gives $\operatorname{reg}(I_{L,K})=4$
over every field. Since
$\operatorname{reg}(I)=\max\{j-i:\beta_{i,j}(I)\ne0\}$, this implies
$\beta_{1,j}(I_{L,K})=0$ for $j\ge6$. This completes the proof over
$\mathbb F_2$.

Finally, orient the graph edges and fundamental cycles over
$\mathbb Z$, and write
$\partial_d:\mathbb Z^{E_d}\to\mathbb Z^{V_d}$ for the oriented
incidence map of $\Gamma_d$. The oriented fundamental cycles form a
$\mathbb Z$-basis of the integral cycle lattice $\ker\partial_d$. The
image of $\partial_d$ is a subgroup of $\mathbb Z^{V_d}$ and hence free,
so $\ker\partial_d$ is a direct summand of $\mathbb Z^{E_d}$ and
$Z_d=K\otimes_{\mathbb Z}\ker\partial_d$ for every field $K$; the
oriented fundamental cycles are thus a $K$-basis of $Z_d$ for every $K$,
and their reductions modulo $2$ are the basis used above. In these
integral bases the same construction gives signed integer matrices for
$\mu_4$ and $\mu_5$ whose reductions modulo $2$ are the two binary
matrices above. Full target
rank modulo $2$ yields a maximal minor with odd determinant, hence a
nonzero maximal minor over $\mathbb Q$. Both multiplication maps
therefore have full target rank over every field of characteristic
zero. Together with the field-independent graph dimensions and the
regularity bound, this proves the characteristic-zero statement.
\end{proof}

\section{Exact verification}

Everything above except the two binary ranks can be checked by hand from
the cover relations of $L$, and those two ranks are recomputable from the
data printed here. One stores the $20$ cover relations $\widehat0<a_i$,
$a_i<b_i$, $a_{i+1}<b_i$, $b_i<\widehat1$, takes the transitive closure
and reads off the $144$ meets and joins; reduces
$x_ax_b-x_{a\wedge b}x_{a\vee b}$ over the $\binom{12}{2}=66$ pairs of
distinct elements, of which $31$ vanish and $35$ give the quadrics
$\mathcal G$ above; forms $\Gamma_d$ for $d=3,4,5$, obtaining $E_d$,
$V_d$, $c_d$ and the fundamental-cycle basis by union--find, so that
$\dim Z_d=E_d-V_d+c_d$; and for $d=4,5$ writes the $12\cdot\dim Z_{d-1}$
products $x_r\zeta$ in that basis and reduces the resulting bit-vector
matrix by Gaussian elimination over $\mathbb F_2$, which returns the
ranks $1528$ and $8625$ of the two matrices displayed above. All
arithmetic is exact, on integers and bit vectors; there is no
randomisation, floating-point computation, external solver or tolerance
parameter, and no step depends on the order of enumeration. The
computation was re-run independently, on a separate machine and by a
program written against this manuscript rather than distributed with it,
and it agreed: all $64$ of its checks passed, covering the lattice, its
purity and rank, its nonmodularity ($40$ failures of the modular law over
all triples $x\le z$, the witness printed above among them), the $35$
generators, every entry of the table, each $\dim Z_d$ recomputed as
$E_d-\operatorname{rank}(\partial_d)$ and certified field-independent by
exact rank bounds, and every row of $\mu_4$ and $\mu_5$ checked against
the cycle space of $\Gamma_d$ over $\mathbb Z$ and over $\mathbb F_2$; it
also goes one degree beyond this note, obtaining $\dim Z_6=35762$ and
$\beta_{1,6}(I_{L,K})=0$ directly. It does not reach
$\beta_{1,j}(I_{L,K})=0$ for $j\ge7$, which no finite computation can
settle; the range $j\ge6$ of Theorem~\ref{thm:main} rests instead on
$\operatorname{reg}(I_{L,K})=4$, quoted from \cite[Lemma~2.8]{HSF} and
not recomputed at all. An independent re-run that agrees is evidence
about the computation, not a second proof. The program used and its
transcript are retained in an auxiliary archive, available on request.

\section{Status and scope}

Theorem~\ref{thm:main} refutes Conjecture~2.4 of \cite{HSF}: the lattice
$L$ is nonmodular and its join--meet ideal is linearly related. Nothing
here bears on the other results of \cite{HSF}, and
\cite[Lemma~2.8]{HSF}, which the proof uses, is taken as correct.

Three limits should be stated. First, $\beta_{1,j}(I_{L,K})=0$ is
obtained by computation only for $j=4,5$; for $j\ge6$ it is deduced from
the quoted regularity bound. Second, the conclusion is established for
$K=\mathbb F_2$ and for $\operatorname{char}K=0$. The dimensions
$\dim_KZ_d$ are field-independent, but full target rank of $\mu_4$ and
$\mu_5$ over a field of characteristic $p\ge3$ is not claimed, and
neither is the value of $\beta_{1,j}(I_{L,K})$ there. Third, no
minimality is claimed for $L$: whether some nonmodular lattice with
fewer than twelve elements has a linearly related join--meet ideal is
not addressed here.

We are aware of no earlier counterexample. A search of the literature on
join--meet ideals and on the linearly related property returned no work
that constructs such a lattice or that otherwise settles
Conjecture~2.4. That search cannot be exhaustive, and the result is
recorded here as refuting the conjecture rather than as closing the
question of priority.

\begin{thebibliography}{9}

\bibitem{HSF}
T.~Hibi and S.~A.~Seyed Fakhari,
\emph{Regularity and depth of binomial ideals arising from combinatorics},
\href{https://arxiv.org/abs/2607.18679v1}{arXiv:2607.18679v1} (2026).

\end{thebibliography}

\end{document}